\unit{Atoms as tests of parity and time-reversal symmetry}
\unitref{Week 12; atomic parity violation and EDM lectures}
\subsection{Why symmetry is a practical experimental tool}
An atomic experiment can be sensitive to
tiny interactions because ordinary
electromagnetic theory predicts that
certain transitions or energy shifts
should be \emph{zero} by symmetry.
If a supposedly forbidden signal
appears in a controlled measurement,
we investigate whether an interaction
absent from the simpler Hamiltonian
has broken that symmetry.
This motivates atomic tests of the
weak interaction, parity violation
and intrinsic electric dipole moments.

\subsection{What parity means}
A parity transformation inverts the
spatial coordinates:
\(\mathbf r\mapsto-\mathbf r\).
For a central potential, the orbital
wavefunctions have definite parity:
\[
P\ket{nlm}=(-1)^l\ket{nlm}.
\]
Thus \(s,d,g\) orbitals are even;
\(p,f\) orbitals are odd.
The electric dipole operator
\(\mathbf d=-e\mathbf r\)
is itself odd under parity.
If initial and final states have
the \emph{same} parity,
\[
\bra f\mathbf d\ket i
=-\bra f\mathbf d\ket i=0
\]
in a parity-conserving model.
This explains the E1 parity selection
rule without memorising a table.

\subsection{How the weak interaction changes the conclusion}
The electromagnetic Hamiltonian of
an atom approximately respects parity,
but the neutral-current weak
interaction contains parity-violating
electron--nucleus terms. Schematically,
\[
H=H_{\rm EM}+H_W,
\quad P H_WP^{-1}\ne H_W.
\]
First-order perturbation theory shows
that an originally definite-parity
state \(\ket i\) acquires admixtures
of opposite-parity states:
\[
\boxed{\ket{\tilde i}\approx\ket i+
\sum_{n\ne i}\ket n
\frac{\bra nH_W\ket i}{E_i-E_n}.}
\]
These admixtures can generate an
otherwise forbidden E1 amplitude
between same-parity levels.

\begin{worked}{Derive a weakly allowed E1 line}
Let \(\ket i,\ket f\) have the same
unperturbed parity and \(\ket n\)
the opposite parity. If only the
initial-state admixture matters,
\[
\ket{\tilde i}=\ket i+\epsilon\ket n,
\quad
\epsilon=\frac{\bra nH_W\ket i}{E_i-E_n}.
\]
Then
\[
\bra f d\ket{\tilde i}
=\underbrace{\bra f d\ket i}_{0}
+\epsilon\bra f d\ket n
=\boxed{\epsilon\bra f d\ket n}.
\]
A small mixing coefficient \(10^{-8}\)
multiplied by an allowed atomic dipole
of order \(ea_0\) gives a
\(10^{-8}ea_0\) amplitude.
In a full calculation, both initial
and final states can be mixed, and
one sums over many intermediate states.
\end{worked}

\subsection{Nuclear weak charge}
One widely discussed contribution
to atomic parity violation is
controlled by a nuclear weak charge,
with tree-level approximate form
\[
\boxed{Q_W\simeq-N+Z(1-4\sin^2\theta_W),}
\]
where \(N,Z\) are neutron and proton
numbers. Electroweak radiative
corrections and nuclear-structure
effects matter in precision work.
The measured parity-violating amplitude
also depends strongly on the atomic
electronic wavefunctions, so comparing
experiment with the Standard Model
requires accurate atomic many-body
theory, not just the weak charge.

\subsection{What is an intrinsic electric dipole moment?}
An electric dipole \(\mathbf d\) changes
sign under parity, while an axial
angular momentum \(\mathbf J\) does not.
A permanent \(\mathbf d\) aligned with
a particle's spin therefore violates
parity. Under time reversal,
spin changes direction but the electric
dipole does not, so such an intrinsic
EDM also violates time reversal.
The interaction energy is
\[
\boxed{H_{\rm EDM}=-\mathbf d\cdot
\boldsymbol{\mathcal E}.}
\]
A spin-precession experiment compares
frequencies for opposite electric-field
directions. A genuine EDM contribution
reverses sign with the electric field,
while carefully chosen combinations
of measurements reject common magnetic
and instrumental backgrounds.

\subsection{Schiff screening and experimental enhancement}
In a simple neutral atom of pointlike,
nonrelativistic Coulomb-interacting
constituents, internal fields screen
a uniformly applied static electric
field at the nucleus. This is the
basis of Schiff's theorem.
Finite nuclear size, relativistic
electron motion and molecular
polarisation circumvent its ideal
assumptions and can make EDM
measurements extremely sensitive.
Heavy atoms and polar molecules
are particularly useful because
their electronic and nuclear
structure can amplify some
symmetry-violating effects.

\subsection{Why this belongs in atomic physics}
The experimental quantity is a tiny
frequency shift or transition amplitude.
To predict it, one still needs the
atomic orbitals, good angular-momentum
coupling, selection rules, perturbation
theory and matrix elements developed
earlier in this mini-text.
A precision experiment turns those
familiar mathematical tools into a
test of fundamental interactions.

\begin{checkpoint}
Parity selection rules forbid amplitudes
\emph{only if} the Hamiltonian respects
parity and the basis states have definite
parity. A weak mixing interaction changes
the state and may open a previously
vanishing transition.
\end{checkpoint}
